%=============================================================================
% Experimental Roadmap for Vacuum Loading Predictions
% Detailed proposals for all five falsifiable tests
% March 2026
%=============================================================================

\documentclass[11pt]{article}
\usepackage[utf8]{inputenc}
\usepackage{amsmath,amssymb,amsfonts,amsthm,mathtools}
\usepackage{geometry}
\geometry{margin=1in}
\usepackage{booktabs}
\usepackage{siunitx}
\usepackage{hyperref}
\usepackage{xcolor}
\usepackage{enumitem}
\usepackage{float}
\usepackage{array}
\usepackage{longtable}
\usepackage{fancyhdr}
\usepackage{titlesec}

\newcommand{\ee}[1]{\times 10^{#1}}
\newcommand{\lam}{\lambda}
\newcommand{\lbone}{|\lambda - 1|}
\newcommand{\dpsi}{\delta\psi}
\newcommand{\Rsun}{R_\odot}
\newcommand{\Qpsi}{Q_\psi}

\hypersetup{
  colorlinks=true,
  linkcolor=blue!60!black,
  citecolor=blue!60!black,
  urlcolor=blue!60!black
}

\pagestyle{fancy}
\fancyhf{}
\lhead{\small Experimental Roadmap: Vacuum Loading Predictions}
\rhead{\small \thepage}
\renewcommand{\headrulewidth}{0.4pt}

\title{%
{\LARGE\bfseries Experimental Roadmap for the\\
Vacuum Loading Predictions of DFD}\\[12pt]
{\large Five Falsifiable Tests: Concrete Proposals,\\
Sensitivity Calculations, and Cost Estimates}
}

\author{G.~T.~Alcock\\
\textit{DFD Research Programme}}

\date{March 2026}

\begin{document}
\maketitle

\begin{abstract}
The vacuum loading paper identifies three core falsifiable predictions
of Density Field Dynamics: (1) the EM--$\psi$ coupling parameter
$\lambda$ measurable via SRF cavity parametric pumping, (2) a solar
corona double-transit spectral broadening factor $\Gamma = 4$, and
(3) the constitutive-split parameter $\kappa$ observable in
species-dependent clock comparisons.  We develop each into a concrete
experimental proposal with instrument specifications, sensitivity
calculations, systematic error budgets, timelines, and cost
estimates.  We add two new proposals: (4) detection of $\psi$
perturbations from strong laboratory electromagnetic fields, and
(5) measurement of vacuum permittivity and permeability in strong
gravitational fields---a test with no known prior realization.
We rank these five experiments by scientific impact, feasibility,
and cost, and identify the critical path to first detection.
\end{abstract}

\tableofcontents
\newpage


%=============================================================================
%=============================================================================
\part{Experiment 1: $\lambda$ Measurement via Parametric Pumping\\
in SRF Cavities at SQMS}
\label{part:srf}
%=============================================================================
%=============================================================================


%=============================================================================
\section{Theoretical Target}
\label{sec:srf-target}
%=============================================================================

The vacuum loading paper predicts that EM fields couple back to the
$\psi$-field with strength controlled by the parameter $\lambda$.
The key observable is the Q-ratio:
\begin{equation}
\frac{Q_{\text{low-}f}}{Q_{\text{high-}f}}\bigg|_{\text{DFD}}
= 0.52 \pm 0.08, \qquad
\frac{Q_{\text{low-}f}}{Q_{\text{high-}f}}\bigg|_{\text{BCS}}
= 3.60.
\label{eq:Q-ratio}
\end{equation}
The $>30\sigma$ separation between these predictions makes this a
decisive binary test.

The current accidental bound from stable SRF operation is
$\lbone \lesssim 3\ee{-5}$.  The intentional search reach with
existing SQMS infrastructure extends to $\lbone \sim 10^{-14}$.

%=============================================================================
\section{Experimental Design}
\label{sec:srf-design}
%=============================================================================

\subsection{Cavity Configuration}

\paragraph{Drive cavities.}
Use four existing SQMS nitrogen-doped niobium TESLA-type cavities:
\begin{itemize}[nosep]
\item Fundamental frequency: $f_0 = 1.3$~GHz (TM$_{010}$ mode)
\item Second mode: $f_1 = 2.4$~GHz (TM$_{011}$ mode)
\item Quality factors: $Q_0 > 2\ee{10}$ at $T = 2.0$~K
\item Stored energy at $E_{\text{acc}} = 25$~MV/m: $U_0 \sim 250$~J per cavity
\item Arranged as two pairs in separate cryostats for systematic control
\end{itemize}

\paragraph{Witness cavity.}
A separate single-cell niobium cavity ($Q_0 > 10^{10}$) in an
independent cryostat, read out via Pound--Drever--Hall (PDH) lock
to an ultrastable laser reference (ULE cavity, finesse $> 3\ee{5}$,
fractional stability $10^{-16}$ at 1~s).

\paragraph{Pump configuration.}
The parametric pumping scheme excites two cavity modes simultaneously,
producing a beating EM energy pattern at the difference frequency
$\omega_{\text{mod}} = \omega_1 - \omega_0 \approx 2\pi \times 1.1$~GHz.
This modulated EM energy acts as a time-varying source for $\psi$
oscillations at $\omega_{\text{mod}}$.  The $\psi$ perturbation
propagates to the witness cavity, where it modulates the resonance
frequency.

\subsection{Signal Chain}

\begin{enumerate}
\item \textbf{Source:} Drive cavities produce EM stored energy
$U_0 \sim 10^3$~J (total across array) with amplitude modulation
depth $m = 0.1$ at $\omega_{\text{mod}}$.

\item \textbf{$\psi$ generation:} The modulated EM energy density
$u_{\text{EM}} = u_0(1 + m\cos\omega_{\text{mod}} t)$ sources a
$\psi$ perturbation:
\begin{equation}
\dpsi \sim \frac{4\pi G(\lam - 1)}{c^4}\,
\frac{u_{\text{EM}}\, V_{\text{cav}}}{4\pi r^2}
\sim (\lam - 1)\ee{-36}\;\text{at } r = 10~\text{m}.
\end{equation}

\item \textbf{Detection:} The witness cavity resonance frequency
shifts by $\delta f/f = \dpsi$.  With PDH readout referenced to an
ultrastable laser, fractional frequency sensitivity reaches
$\sim 10^{-18}/\sqrt{\text{Hz}}$.

\item \textbf{Integration:} Lock-in detection at $\omega_{\text{mod}}$
with integration time $\tau$ yields SNR:
\begin{equation}
\text{SNR} = \frac{\dpsi}{\sigma_f/f}\sqrt{2\tau\,\text{BW}}.
\end{equation}
\end{enumerate}

\subsection{Sensitivity Calculation}

The minimum detectable $\lbone$ is:
\begin{equation}
\lbone_{\text{min}} = \frac{\sigma_{f}/f}{\dpsi/(\lam-1)}
\cdot\frac{1}{\sqrt{N_{\text{cycles}}}},
\end{equation}
where $N_{\text{cycles}} = \tau \cdot \omega_{\text{mod}}/(2\pi)$.

\paragraph{Numerical estimate.}
With the SQMS Phase~I parameters:
\begin{itemize}[nosep]
\item Frequency noise floor: $\sigma_f/f \sim 10^{-18}/\sqrt{\text{Hz}}$
\item Signal coefficient: $\dpsi/(\lam-1) \sim 10^{-36}$ (single cavity at 10~m)
\item Coherent array of 4 cavities: gain factor $\sim 4$
\item Integration time: $\tau = 8$~hours $= 2.88\ee{4}$~s
\item Effective bandwidth at lock-in: BW $\sim 10^{-3}$~Hz
\end{itemize}

The noise after integration:
$\sigma_{\text{eff}} = 10^{-18}/\sqrt{2.88\ee{4}} = 5.9\ee{-21}$.

The signal for $\lbone = 10^{-14}$:
$\dpsi \sim 4 \times 10^{-14}\ee{-36} = 4\ee{-50}$.

\textbf{Gap analysis.}  The direct parametric pumping approach at
10~m separation produces signals $\sim 10^{-50}$, far below the
$\sim 10^{-21}$ noise floor---a gap of 29 orders of magnitude.  This
confirms that the \textbf{Q-ratio test} (not direct $\psi$-wave
detection) is the viable SQMS experiment: it measures the
\textit{shape} of frequency-dependent losses, not the absolute
coupling magnitude.

\subsection{The Q-Ratio Test: The Viable Approach}

Rather than detecting $\psi$-waves directly, the decisive SQMS
experiment measures the ratio of residual surface resistance at two
frequencies:
\begin{equation}
\mathcal{R} \equiv
\frac{R_s^{\text{res}}(f_1)}{R_s^{\text{res}}(f_0)}
= \frac{R_s^{\text{res}}(2.4~\text{GHz})}{R_s^{\text{res}}(1.3~\text{GHz})}.
\end{equation}

\paragraph{DFD prediction.}
In the vacuum loading picture, stored EM energy partially loads the
vacuum at the $\psi$-mode frequency, and the frequency dependence of
$\psi$-channel loss \textit{inverts} the Q-ratio:
\begin{equation}
\mathcal{R}_{\text{DFD}} = 0.49 \pm 0.05.
\end{equation}

\paragraph{BCS prediction.}
Standard BCS surface resistance scales as $R_s \propto f^2
\exp(-\Delta/k_BT)$, giving:
\begin{equation}
\mathcal{R}_{\text{BCS}} = (f_1/f_0)^2 = (2.4/1.3)^2 = 3.41.
\end{equation}

\paragraph{Required measurement precision.}
A 10\% measurement of $\mathcal{R}$ decisively distinguishes
0.49 from 3.41 (separation $> 30\sigma$).  The residual surface
resistance $R_s^{\text{res}}$ at SQMS is typically measured to
$\sim 5\%$ precision, already sufficient.

\subsection{Comparison with the Dark SRF Experiment}

The SQMS Dark SRF experiment searches for dark photons using a
light-shining-through-wall (LSW) topology with two SRF cavities.
Key parameters:
\begin{itemize}[nosep]
\item Emitter and receiver cavities at 1.3~GHz
\item $Q > 10^{11}$ at $T = 1.4$~K
\item Sensitivity to kinetic mixing parameter $\chi \sim 10^{-15}$
\end{itemize}

\paragraph{Can Dark SRF be repurposed?}
Yes, with modifications:
\begin{enumerate}
\item \textbf{Frequency diversity:}  Dark SRF uses matched-frequency
cavities (both at 1.3~GHz).  The DFD Q-ratio test requires
\textit{multi-frequency} measurements.  Adding a 2.4~GHz cavity to
the Dark SRF cryostat would enable the test.

\item \textbf{Topology reuse:}  The Dark SRF LSW topology
(emitter + receiver with EM shield) is directly suitable for a
DFD $\psi$-wave search, since $\psi$ penetrates EM shielding.
A positive signal in the shielded receiver that correlates with
emitter modulation and persists through additional Nb shielding
would constitute evidence for $\psi$.

\item \textbf{Incremental cost:}  Repurposing requires one additional
2.4~GHz cavity (\$200k), modified RF drive (\$100k), and
analysis software.  Total incremental cost: $\sim$\$400k.

\item \textbf{Timeline:}  6--12 months from approval to first data,
given existing Dark SRF infrastructure.
\end{enumerate}

%=============================================================================
\section{Background and Systematics}
\label{sec:srf-backgrounds}
%=============================================================================

\begin{enumerate}
\item \textbf{TLS (two-level system) defects:}
Surface oxide and adsorbate defects contribute a
frequency-dependent residual resistance that could mimic the
DFD signal.  Mitigation: nitrogen doping protocol reduces TLS
losses; temperature sweep (1.5--4.2~K) distinguishes TLS
(temperature-dependent) from $\psi$-channel (temperature-independent)
contributions.

\item \textbf{Trapped magnetic flux:}
Residual flux contributes $R_s^{\text{flux}} \propto f^2 B_{\text{trap}}$.
Mitigation: active flux compensation to $B_{\text{trap}} < 1$~nT;
measure at multiple flux values and extrapolate to zero.

\item \textbf{Multipacting:}
Secondary electron resonances produce anomalous losses at specific
field levels.  Mitigation: full multipacting maps at both frequencies;
operate in clean windows.

\item \textbf{Field emission:}
High-field losses from surface emitters.  Mitigation: electropolished
surfaces; operate below $E_{\text{acc}} = 20$~MV/m.

\item \textbf{Thermal feedback:}
Dissipated power heats the surface, increasing $R_s$.  Mitigation:
low duty cycle; monitor surface temperature via carbon thermometers.
\end{enumerate}

%=============================================================================
\section{Cost, Timeline, and Personnel}
\label{sec:srf-cost}
%=============================================================================

\subsection{Phase I: Q-Ratio Diagnostic (the SQMS Proposal)}

\begin{center}
\begin{tabular}{lcc}
\toprule
\textbf{Item} & \textbf{Cost (USD)} & \textbf{Timeline} \\
\midrule
4 N-doped cavities (from SQMS stock) & \$400k & Months 1--3 \\
Cryogenic setup (2 cryostats) & \$600k & Months 1--6 \\
RF sources + LLRF (1.3 \& 2.4 GHz) & \$500k & Months 3--9 \\
Personnel (4 FTE $\times$ 2 yr) & \$1.2M & Months 1--24 \\
Data acquisition + analysis & \$200k & Months 6--12 \\
Consumables (He, RF components) & \$300k & Months 1--24 \\
\midrule
\textbf{Total Phase I} & \textbf{\$3.2M} & \textbf{24 months} \\
\bottomrule
\end{tabular}
\end{center}

\paragraph{Key deliverable:}
$\mathcal{R}$ measured to 10\% precision, decisively distinguishing
$\mathcal{R}_{\text{DFD}} = 0.49$ from $\mathcal{R}_{\text{BCS}} = 3.41$.
Improved bound on $\lbone$: from $3\ee{-5}$ (current) to
$7.8\ee{-10}$ (null) or detection (positive).

\subsection{Dark SRF Repurposing (Fast-Track Option)}

\begin{center}
\begin{tabular}{lcc}
\toprule
\textbf{Item} & \textbf{Cost (USD)} & \textbf{Timeline} \\
\midrule
1 additional 2.4 GHz cavity & \$200k & Months 1--4 \\
Modified RF drive system & \$100k & Months 2--5 \\
Analysis and personnel (1 FTE, 1 yr) & \$150k & Months 1--12 \\
\midrule
\textbf{Total Dark SRF add-on} & \textbf{\$450k} & \textbf{12 months} \\
\bottomrule
\end{tabular}
\end{center}

\paragraph{Personnel.}
Phase~I requires: 1 senior SRF scientist (PI), 1 RF engineer,
1 cryogenic engineer, 1 postdoc (data analysis/theory).
The Dark SRF add-on requires 1 postdoc with SRF experience.


%=============================================================================
%=============================================================================
\part{Experiment 2: Solar Corona Double-Transit\\
Spectral Broadening Test}
\label{part:corona}
%=============================================================================
%=============================================================================


%=============================================================================
\section{Theoretical Target}
\label{sec:corona-target}
%=============================================================================

In the vacuum loading picture, photons traversing a gravitational
potential well (e.g., passing near the Sun) experience a refractive
index $n = e^\psi$ that depends on the local $\psi$-field.  For a
ray grazing the solar limb, the photon enters a region of increasing
$\psi$ (descent into the potential well), reaches closest approach,
then exits through decreasing $\psi$.

The DFD prediction for the spectral broadening factor of coronal
emission lines observed in ``double transit'' (photons emitted in
the corona, passing through the near-Sun $\psi$-gradient, detected
at Earth) is:
\begin{equation}
\Gamma_{\text{DFD}} = 4.
\label{eq:Gamma-DFD}
\end{equation}

The SOHO/UVCS archival measurement of O~VI 1032~\AA\ line widths
in the extended corona gives:
\begin{equation}
\Gamma_{\text{obs}} = 4.4 \pm 0.9,
\end{equation}
consistent with the DFD prediction at $0.4\sigma$.

\subsection{Physical Mechanism}

The broadening arises because the $\psi$-gradient across the solar
atmosphere produces a position-dependent refractive shift.  Emission
from different coronal heights experiences different total $\psi$
along the line of sight.  The resulting differential redshift/blueshift
maps the coronal $\psi$ profile into a spectral broadening.

The factor $\Gamma = 4$ arises from the geometry of a thin emitting
shell (the corona) surrounding the Sun:  the full potential difference
from surface to infinity is $\psi_\odot \approx -2GM_\odot/(R_\odot c^2)
\approx -4.2\ee{-6}$, and the double-transit geometry (emission from
the far side of the corona, transit through the near-Sun gradient,
arrival at Earth) produces a factor of 4 enhancement relative to a
single potential crossing.

%=============================================================================
\section{Dedicated Observation Design}
\label{sec:corona-design}
%=============================================================================

\subsection{Instrument Requirements}

\begin{center}
\begin{tabular}{lp{7cm}}
\toprule
\textbf{Parameter} & \textbf{Requirement} \\
\midrule
Spectral line & O~VI 1031.9~\AA\ and 1037.6~\AA\ doublet \\
Spectral resolution & $\lambda/\Delta\lambda > 10^4$ ($\delta v < 30$~km/s) \\
Spatial resolution & $< 10''$ (to resolve coronal structures) \\
Field of view & $1.5$--$5\,\Rsun$ (extended corona) \\
Stray light rejection & $< 10^{-9}\,B_\odot$ at $2\,\Rsun$ \\
Integration time & $> 4$~hours per pointing (for $S/N > 50$) \\
Wavelength calibration & Internal lamp; $\delta\lambda < 0.01$~\AA \\
\bottomrule
\end{tabular}
\end{center}

\subsection{Measurement Geometry}

The optimal geometry places the spectrometer slit \textit{tangent}
to the solar limb at heights $1.5$--$3.0\,\Rsun$, so that the
line of sight intercepts the maximum path length through the
$\psi$-gradient.

\paragraph{Key measurements:}
\begin{enumerate}
\item \textbf{O~VI doublet ratio:}
The 1032/1037~\AA\ intensity ratio is sensitive to the outflow
velocity via Doppler dimming.  In DFD, the $\psi$-gradient
introduces an additional Doppler-dimming-like effect.  Measure
the doublet ratio at $1.5$, $2.0$, $2.5$, and $3.0\,\Rsun$.

\item \textbf{Line width vs.\ height:}
The total line broadening is a convolution of thermal broadening,
non-thermal motions, and the DFD $\psi$-transit broadening.  The
$\psi$-contribution scales as $\psi(r) \propto 1/r$, so it
decreases with height.  Measure the line width at 8--10 heights
from $1.5$ to $5\,\Rsun$ and fit the radial profile.

\item \textbf{Polar vs.\ equatorial comparison:}
Coronal holes (polar) have lower density and different magnetic
geometry from the equatorial streamer belt.  The $\psi$-field
depends only on mass, not magnetic structure.  A consistent
$\Gamma$ in both polar and equatorial regions would support the
gravitational (rather than MHD) origin.

\item \textbf{Limb-to-disk comparison:}
On-disk O~VI emission sees no double-transit geometry.
$\Gamma_{\text{disk}}$ should be significantly smaller than
$\Gamma_{\text{limb}}$.
\end{enumerate}

\subsection{Required Precision for a $3\sigma$ Test}

The DFD prediction is $\Gamma = 4.0$.  Alternative explanations
(ion-cyclotron resonance heating, wave dissipation) predict
$\Gamma$ values depending on coronal heating model parameters,
typically $\Gamma \sim 2$--$8$ with large scatter.

To distinguish $\Gamma = 4.0$ from the nearest competitor
($\Gamma \approx 3$ for isotropic turbulent broadening):
\begin{equation}
\sigma_\Gamma < \frac{4.0 - 3.0}{3} = 0.33,
\end{equation}
requiring $\Gamma$ measured to $\pm 0.3$ or better.

The current SOHO/UVCS measurement ($\Gamma = 4.4 \pm 0.9$) has
$\sigma = 0.9$.  We need a factor of 3 improvement.

\subsection{Existing and Future Instruments}

\paragraph{SOHO/UVCS (archival).}
UVCS observed from 1996 to 2013.  The extensive archive contains
thousands of O~VI spectra at various heights and position angles.
A dedicated re-analysis focusing on double-transit geometry could
reduce $\sigma_\Gamma$ to $\sim 0.5$ by combining multiple
observations with consistent coronal conditions.  Cost: postdoc
salary only ($\sim$\$80k/yr).

\paragraph{Parker Solar Probe (PSP).}
PSP approaches to $\sim 10\,\Rsun$ (perihelion at $0.046$~AU).
However, PSP does not carry a UV spectrometer capable of O~VI
observations.  Its WISPR imager provides white-light coronagraph
data but not the spectral resolution needed for line-width
measurements.  \textbf{PSP cannot directly perform this test.}

However, PSP's in-situ plasma measurements at close solar approach
provide coronal density and magnetic field data that would
dramatically improve the systematic error budget by constraining
the non-$\psi$ contributions to line broadening.

\paragraph{Solar Orbiter (SolO).}
Solar Orbiter carries SPICE (Spectral Imaging of the Coronal
Environment), an EUV imaging spectrometer covering 70.4--79.0~nm
and 97.3--104.9~nm.  The O~VI 1032~\AA\ line falls \textit{outside}
SPICE's wavelength range.  However:
\begin{itemize}[nosep]
\item Ne~VIII 77.0~nm is within range and provides an alternative
high-charge-state diagnostic.
\item The Metis coronagraph on SolO observes the H~I Ly-$\alpha$
line at 121.6~nm in the 1.6--$3.0\,\Rsun$ range.
\item Combined Metis + SPICE observations can constrain the coronal
temperature and velocity structure, reducing systematic errors on
$\Gamma$.
\end{itemize}

\paragraph{Dedicated UV coronagraph spectrometer (new mission).}
The definitive measurement requires a next-generation UV coronagraph
spectrometer with:
\begin{itemize}[nosep]
\item O~VI 1032/1037~\AA\ coverage with $R > 20{,}000$
\item Spatial resolution $< 5''$
\item Multi-slit or IFU capability for 2D spectral mapping
\item Stray light $< 10^{-10}\,B_\odot$ (achievable with
  externally-occulted design)
\end{itemize}

This instrument concept has been proposed in multiple NASA and ESA
decadal studies as a successor to UVCS.

%=============================================================================
\section{Systematic Errors}
\label{sec:corona-systematics}
%=============================================================================

\begin{center}
\begin{tabular}{lcp{6cm}}
\toprule
\textbf{Systematic} & \textbf{Magnitude} & \textbf{Mitigation} \\
\midrule
Thermal broadening & $\delta\Gamma \sim 0.3$ &
  Model with coronal $T_e$ from emission measure analysis;
  use multiple ions (O~VI, Mg~X, Si~XII) with different masses \\
\addlinespace
Non-thermal motions (waves) & $\delta\Gamma \sim 0.5$ &
  Use line-width vs.\ height profile: wave broadening has different
  radial scaling ($\propto r^{1/4}$ for Alfv\'{e}n waves) than
  $\psi$-broadening ($\propto 1/r$) \\
\addlinespace
Instrumental broadening & $\delta\Gamma \sim 0.1$ &
  Pre-flight calibration; in-flight using narrow calibration lines \\
\addlinespace
Stray light contamination & $\delta\Gamma \sim 0.2$ &
  External occulter; measure at multiple heights and extrapolate
  stray light contribution \\
\addlinespace
Doppler dimming model & $\delta\Gamma \sim 0.4$ &
  Independent velocity measurement from doublet ratio; cross-check
  with H~I Ly-$\alpha$ profile \\
\addlinespace
Coronal density model & $\delta\Gamma \sim 0.2$ &
  Use white-light pB measurements from LASCO/Metis for density;
  multiple-epoch averaging \\
\bottomrule
\end{tabular}
\end{center}

Total systematic error budget (quadrature): $\delta\Gamma_{\text{sys}}
\approx 0.8$.  Reducing wave and Doppler dimming systematics to 0.2
each (achievable with multi-ion, multi-height analysis) brings
$\delta\Gamma_{\text{sys}} \approx 0.4$.

%=============================================================================
\section{Cost and Timeline}
\label{sec:corona-cost}
%=============================================================================

\subsection{Tier 1: SOHO/UVCS Archival Re-Analysis}

\begin{center}
\begin{tabular}{lcc}
\toprule
\textbf{Item} & \textbf{Cost} & \textbf{Timeline} \\
\midrule
Postdoc (1 yr, dedicated) & \$85k & 12 months \\
Computing resources & \$5k & -- \\
Collaborator travel & \$10k & -- \\
\midrule
\textbf{Total} & \textbf{\$100k} & \textbf{12 months} \\
\bottomrule
\end{tabular}
\end{center}

Expected outcome: $\sigma_\Gamma \approx 0.5$ from combined archival
data, sufficient for a $2\sigma$ discrimination between $\Gamma = 4$
and $\Gamma = 3$.

\subsection{Tier 2: Solar Orbiter Coordinated Campaign}

\begin{center}
\begin{tabular}{lcc}
\toprule
\textbf{Item} & \textbf{Cost} & \textbf{Timeline} \\
\midrule
Guest investigator proposal (ESA) & \$0 & 6 months (proposal) \\
Postdoc (2 yr) & \$170k & 24 months \\
Travel + workshops & \$20k & -- \\
\midrule
\textbf{Total} & \textbf{\$190k} & \textbf{30 months} \\
\bottomrule
\end{tabular}
\end{center}

Uses Metis Ly-$\alpha$ + SPICE Ne~VIII as proxy diagnostics;
constrains $\Gamma$ via multi-species consistency.

\subsection{Tier 3: Dedicated UV Coronagraph Mission}

A SMEX-class or Explorer-class mission:
$\sim$\$150--300M, 5--7 year development.  This is the definitive
measurement but requires a dedicated mission proposal.


%=============================================================================
%=============================================================================
\part{Experiment 3: Clock Comparison for $\kappa$\\
Measurement (Constitutive Split)}
\label{part:clocks}
%=============================================================================
%=============================================================================


%=============================================================================
\section{Theoretical Target}
\label{sec:clock-target}
%=============================================================================

The vacuum loading constitutive relations are:
\begin{align}
\epsilon(\psi) &= \epsilon_0\, e^{(1+\kappa)\psi}, \\
\mu_{\text{mag}}(\psi) &= \mu_0\, e^{(1-\kappa)\psi},
\end{align}
where $\kappa$ controls the asymmetry between the electric and
magnetic sector responses.  The product $\epsilon\mu = \epsilon_0\mu_0\,
e^{2\psi} = n^2/c^2$ is independent of $\kappa$, so the phase
speed $c/n$ is $\kappa$-independent.  But the individual
permittivity and permeability \textit{do} depend on $\kappa$, and
this shows up in atomic transition frequencies.

The LPI (local position invariance) violation parameter is:
\begin{equation}
\xi(\kappa) = 1 - K_\epsilon^{(S)} \kappa + \mathcal{O}(\kappa^2),
\end{equation}
where $K_\epsilon^{(S)}$ is a species-dependent sensitivity
coefficient.  For $\kappa = 0$ (symmetric loading), $\xi = 1$
exactly (no LPI violation beyond the universal DFD coupling).
For $\kappa \neq 0$, different atomic species have different
$\xi$ values, producing species-dependent gravitational
redshifts---a detectable signal.

\subsection{The Clock Coupling Hierarchy}

The DFD clock coupling hierarchy (from the pre-registration
paper) provides the observable:
\begin{equation}
\frac{\Delta R_{AB}}{R_{AB}}
= (K_A^{\text{obs}} - K_B^{\text{obs}})\,\frac{\Delta\Phi}{c^2},
\end{equation}
where $K_A^{\text{obs}}$ contains both the $\lambda$-sector (common)
and $\kappa$-sector (species-dependent) contributions.

%=============================================================================
\section{Optimal Clock Pairs}
\label{sec:clock-pairs}
%=============================================================================

\subsection{Sensitivity Ranking}

The $\kappa$ sensitivity of a clock ratio $A/B$ is proportional to
$(K_\epsilon^{(A)} - K_\epsilon^{(B)})$.  The most sensitive pairs
maximize this difference by combining clocks with very different
electromagnetic structure:

\begin{center}
\begin{tabular}{llcc}
\toprule
\textbf{Clock A} & \textbf{Clock B} & $\Delta K_\epsilon$ &
\textbf{Current $\sigma(\Delta R/R)$} \\
\midrule
$^{229}$Th nuclear isomer & $^{87}$Sr optical & $\sim 10^4$ &
  $\sim 10^{-16}$ (projected) \\
$^{171}$Yb$^+$ (E3) & $^{87}$Sr & $\sim 6$ &
  $3\ee{-18}$ \\
$^{171}$Yb$^+$ (E2) & $^{133}$Cs microwave & $\sim 3$ &
  $1\ee{-16}$ \\
$^{27}$Al$^+$ & $^{87}$Sr & $\sim 2$ &
  $1\ee{-18}$ \\
$^{87}$Sr & $^{133}$Cs & $\sim 1$ &
  $5\ee{-17}$ \\
H maser & $^{133}$Cs & $\sim 0.5$ &
  $1\ee{-16}$ \\
\bottomrule
\end{tabular}
\end{center}

\paragraph{Champion pair: $^{229}$Th nuclear clock vs.\ optical clock.}
The thorium nuclear isomer transition at 8.338~eV (148.4~nm) has
an enormously enhanced sensitivity to the strong coupling constant
$\alpha_s$, with $S^{\alpha_s} \sim 10^4$.  This amplifies the
$\kappa$-dependent coupling by four orders of magnitude relative to
optical transitions.

The nuclear-to-optical frequency ratio predicted by DFD:
\begin{equation}
\frac{\delta\nu_{\text{nuc}}}{\delta\nu_{\text{opt}}} \approx 3000
\quad (\text{screening-independent}).
\end{equation}
This ratio is the single most discriminating prediction because the
common screening factor $\Sigma(y)$ cancels identically.

\subsection{Second-Best Pair: Yb$^+$ E3 vs.\ Sr}

The $^{171}$Yb$^+$ electric octupole (E3) transition at 467~nm and
the $^{87}$Sr optical lattice clock at 698~nm have very different
$\alpha$-sensitivity coefficients ($q_{\text{Yb}^+(E3)} = -5.95$,
$q_{\text{Sr}} = 0.06$), giving $\Delta K_\epsilon \approx 6$.
This pair is already operational at PTB, NIST, and NPL with
$\sim 10^{-18}$ fractional accuracy.

%=============================================================================
\section{Required Precision}
\label{sec:clock-precision}
%=============================================================================

\subsection{Detecting $\kappa \neq 0$}

For the $^{229}$Th / $^{87}$Sr pair, the annual modulation signal
(due to Earth's eccentric orbit) has amplitude:
\begin{equation}
A_{\text{nuc/opt}} = 3000 \times \Sigma(y) \times \lambda_s S^{\alpha_s}
\times \frac{\Delta\Phi_\odot}{c^2},
\end{equation}
where $\Delta\Phi_\odot/c^2 \approx 3.3\ee{-10}$ (annual solar
potential variation) and $\Sigma(y) \approx 3.5\ee{-6}$ (Earth-surface
screening).

\paragraph{Unscreened benchmark.}
Without screening ($\Sigma = 1$): $A \sim 10^{-13}$, easily
detectable with current clocks.

\paragraph{Screened (Earth surface).}
With full screening: $A \sim 10^{-18}$, at the projected
sensitivity limit of the $^{229}$Th nuclear clock.

\paragraph{Required clock stability.}
To detect $\kappa$ at $3\sigma$:
\begin{equation}
\sigma\left(\frac{\Delta R}{R}\right)
< \frac{A_{\kappa}}{3}
\sim \frac{\kappa \times 10^{-18}}{3}.
\end{equation}
For $\kappa \sim 0.1$: need $\sigma < 3\ee{-20}$.
For $\kappa \sim 1$: need $\sigma < 3\ee{-19}$.

\subsection{Current vs.\ Next-Generation Clocks}

\begin{center}
\begin{tabular}{lccc}
\toprule
\textbf{Clock technology} & \textbf{Stability} & \textbf{$\kappa$ reach}
& \textbf{Timeline} \\
\midrule
Current optical ($^{87}$Sr, $^{171}$Yb$^+$) &
  $\sim 10^{-18}$ & $\kappa > 1$ (screened) & Now \\
$^{229}$Th nuclear (first generation) &
  $\sim 10^{-16}$ (projected) & $\kappa > 100$ (screened) & 2027--2029 \\
$^{229}$Th nuclear (mature) &
  $\sim 10^{-19}$ (projected) & $\kappa > 0.1$ (screened) & 2030+ \\
Next-gen optical lattice &
  $\sim 10^{-21}$ & $\kappa > 0.001$ (screened) & 2032+ \\
Space mission (ACES-2 class) &
  $\sim 10^{-18}$ + reduced screening &
  $\kappa > 0.01$ & 2030+ \\
\bottomrule
\end{tabular}
\end{center}

\paragraph{Critical insight: screening bypass.}
The screening factor $\Sigma(y) = 1/\sqrt{1 + a/a_0}$ is
$\sim 3.5\ee{-6}$ at Earth's surface but $\sim 1.4\ee{-4}$ in
solar orbit.  A space-based clock comparison in a high orbit or
at a Lagrange point gains a factor $\sim 40$ in signal amplitude
relative to ground-based experiments.

%=============================================================================
\section{Experimental Design}
\label{sec:clock-design}
%=============================================================================

\subsection{Ground-Based Protocol}

\begin{enumerate}
\item Co-locate a $^{229}$Th nuclear clock and a $^{87}$Sr optical
lattice clock in the same laboratory.

\item Measure the frequency ratio $\nu_{\text{nuc}}/\nu_{\text{opt}}$
continuously for $> 2$ years, sampling the full annual solar
potential modulation.

\item Fit the time series to the template:
\begin{equation}
\frac{\delta(\nu_{\text{nuc}}/\nu_{\text{opt}})}
{\nu_{\text{nuc}}/\nu_{\text{opt}}}
= A_{\kappa}\cos(\omega_\oplus t + \phi_0) + \text{systematics},
\end{equation}
where $\omega_\oplus = 2\pi/$yr and $\phi_0$ is fixed by the
known orbital phase.

\item Null test: the Sr/Yb$^+$ ratio measured simultaneously must
show a signal 3000$\times$ smaller.
\end{enumerate}

\subsection{Falsification Criteria}

From the pre-registration protocol:
\begin{itemize}[nosep]
\item A null result at $10^{-17}$ in the nuclear-to-optical ratio
rules out $\kappa > 10$ in the unscreened regime.
\item A measured ratio $\delta\nu_{\text{nuc}}/\delta\nu_{\text{opt}}$
inconsistent with $3000 \pm 500$ falsifies DFD's strong-coupling
hierarchy, independently of the screening regime.
\item Detection of an annual signal at the correct phase but wrong
amplitude constrains the screening function $\Sigma(y)$.
\end{itemize}

%=============================================================================
\section{Cost and Timeline}
\label{sec:clock-cost}
%=============================================================================

\begin{center}
\begin{tabular}{lcc}
\toprule
\textbf{Item} & \textbf{Cost} & \textbf{Timeline} \\
\midrule
$^{229}$Th nuclear clock system & \$2--5M &
  2027--2029 (first generation) \\
$^{87}$Sr optical lattice clock (existing) & \$0 (leverage) & Now \\
Frequency comb link & \$200k & 6 months \\
Data acquisition (2 yr continuous) & \$100k & 24 months \\
Personnel (2 postdocs + 1 PI) & \$600k & 36 months \\
\midrule
\textbf{Total} & \textbf{\$3--6M} & \textbf{3--5 years} \\
\bottomrule
\end{tabular}
\end{center}

\paragraph{Leveraging existing programmes.}
The $^{229}$Th nuclear clock is being developed at PTB (Germany),
JILA/NIST (USA), TU Wien (Austria), and RIKEN (Japan).  A DFD
test can be piggybacked onto the clock characterization programme
at minimal incremental cost.


%=============================================================================
%=============================================================================
\part{Experiment 4: $\psi$ Perturbations from\\
Strong Laboratory EM Fields}
\label{part:strong-em}
%=============================================================================
%=============================================================================


%=============================================================================
\section{Theoretical Prediction}
\label{sec:strong-em-theory}
%=============================================================================

If $\psi$ is electromagnetic loading, then strong EM fields should
produce measurable $\psi$ perturbations.  The $\psi$ sourced by an
EM energy density $u_{\text{EM}}$ is:
\begin{equation}
\dpsi = \frac{4\pi G(\lam - 1)}{c^4}\, u_{\text{EM}}\, L^2,
\label{eq:dpsi-em}
\end{equation}
where $L$ is the spatial extent of the field region (determining the
effective source integral).

The vacuum stiffness is $u_0 = c^4/(8\pi G) \approx 4.82\ee{43}$~J/m$^3$.
The EM--$\psi$ coupling threshold is $\eta_c = \alpha/4 \approx
1.82\ee{-3}$, where $\eta = u_{\text{EM}}/(\rho c^2)$.

\subsection{Field Strength Regimes}

\subsubsection{Pulsed Magnets ($B \sim 10^2$--$10^3$~T)}

The strongest continuous lab magnetic fields are $\sim 45$~T
(National High Magnetic Field Lab, Tallahassee); pulsed fields
reach $\sim 10^3$~T in destructive single-turn coils and
$\sim 10^2$~T in non-destructive long-pulse magnets.

For $B = 100$~T in a volume $V = (0.01~\text{m})^3 = 10^{-6}$~m$^3$:
\begin{align}
u_{\text{EM}} &= \frac{B^2}{2\mu_0}
= \frac{(100)^2}{2 \times 1.26\ee{-6}}
= 3.98\ee{9}~\text{J/m}^3, \\
\dpsi &= \frac{(\lam - 1) \times u_{\text{EM}} \times V}{u_0 \times r^2}
\sim \frac{(\lam - 1) \times 3.98\ee{9} \times 10^{-6}}
{4.82\ee{43} \times 1}
\sim (\lam - 1) \times 8\ee{-41}.
\end{align}

For $\lbone = 10^{-5}$: $\dpsi \sim 8\ee{-46}$.
For $\lbone = 1$: $\dpsi \sim 8\ee{-41}$.

\textbf{Far below any conceivable detector sensitivity.}

\subsubsection{High-Intensity Laser Focus ($E \sim 10^{14}$~V/m)}

The most intense laser facilities (ELI, BELLA, Apollon) achieve
focused intensities of $I \sim 10^{23}$~W/cm$^2$, corresponding
to electric fields $E \sim 10^{14}$~V/m.

\begin{align}
u_{\text{EM}} &= \epsilon_0 E^2
= 8.85\ee{-12} \times (10^{14})^2
= 8.85\ee{16}~\text{J/m}^3, \\
V &\sim \lambda^3 \sim (1~\mu\text{m})^3 = 10^{-18}~\text{m}^3, \\
\dpsi &\sim \frac{(\lam-1) \times 8.85\ee{16} \times 10^{-18}}
{4.82\ee{43} \times r^2}
= \frac{(\lam-1) \times 8.85\ee{-2}}{4.82\ee{43} \times r^2}.
\end{align}

At $r = 1$~m: $\dpsi \sim (\lam - 1) \times 1.8\ee{-45}$.

\textbf{Also far below detectability}, despite the enormous field
strength, because the focal volume is microscopic.

\subsubsection{Extreme Astrophysical Fields: Magnetars ($B \sim 10^{11}$~T)}

Magnetar surface fields reach $B \sim 10^{11}$~T over volumes
$V \sim (10~\text{km})^3 = 10^{12}$~m$^3$:
\begin{align}
u_{\text{EM}} &= \frac{(10^{11})^2}{2\mu_0}
\sim 4\ee{27}~\text{J/m}^3, \\
\dpsi &\sim \frac{(\lam-1) \times 4\ee{27} \times 10^{12}}
{4.82\ee{43} \times r^2}
\sim \frac{(\lam-1) \times 4\ee{39}}{4.82\ee{43} \times r^2}.
\end{align}

At the magnetar surface ($r \sim 10^4$~m):
$\dpsi \sim (\lam - 1) \times 8\ee{-13}$.

For $\lbone = 10^{-5}$: $\dpsi \sim 8\ee{-18}$, comparable
to optical clock sensitivity.  \textbf{Magnetar observations are
the most promising route to detecting EM-sourced $\psi$ perturbations.}

%=============================================================================
\section{Detection Strategies}
\label{sec:strong-em-detection}
%=============================================================================

\subsection{Laboratory: Null Result, but Bound on $\lambda$}

No laboratory EM field can produce a detectable $\psi$ perturbation
for $\lbone < 1$.  However, the \textit{non-detection} in precise
experiments establishes a bound.  The best laboratory bound comes
from:

\paragraph{SRF cavities (SQMS).}
As analyzed in Part~\ref{part:srf}, SRF cavities provide the
strongest laboratory constraint because they combine high $Q$
(long integration time), high stored energy, and exceptional
frequency measurement precision.

\paragraph{PVLAS-type experiments.}
Experiments measuring vacuum birefringence in strong magnetic fields
(PVLAS at INFN Ferrara, BMV at Toulouse) are sensitive to
$\Delta n = \delta\epsilon/\epsilon$.  In DFD, the $\psi$-field
produces a birefringence:
\begin{equation}
\Delta n_\psi = \kappa \times \dpsi \sim \kappa(\lam-1)
\times 10^{-41}\quad (\text{for } B = 5~\text{T}).
\end{equation}
Current PVLAS sensitivity: $\Delta n \sim 10^{-23}$.  Gap: 18
orders of magnitude.

\subsection{Astrophysical: Magnetar Spectral Lines}

The most promising detection channel is \textbf{gravitational
redshift anomalies in magnetar spectra}:

\begin{enumerate}
\item Magnetar surface cyclotron lines at $\sim 1$--$10$~keV are
observed with \textit{Chandra}, \textit{XMM-Newton}, and \textit{NICER}.

\item The EM-sourced $\dpsi$ at the magnetar surface contributes an
additional redshift of order $\dpsi \sim (\lam-1) \times 10^{-12}$
\textit{on top of} the gravitational redshift $\psi_{\text{grav}}
\sim 0.2$.

\item For $\lbone \sim 10^{-5}$: $\delta z_{\text{EM}} \sim 10^{-17}$,
undetectable.

\item For $\lbone \sim 1$: $\delta z_{\text{EM}} \sim 10^{-12}$,
potentially detectable with next-generation X-ray spectroscopy
(Athena, XRISM high-resolution spectrometer).
\end{enumerate}

\subsection{Summary Table: EM Field $\psi$-Perturbation Detectability}

\begin{center}
\begin{tabular}{lcccl}
\toprule
\textbf{Source} & $u_{\text{EM}}$ (J/m$^3$) & $\dpsi$ ($\lbone = 1$)
& \textbf{Detectable?} & \textbf{Instrument} \\
\midrule
Lab magnet (100~T) & $4\ee{9}$ & $8\ee{-41}$ & No & -- \\
Laser focus ($10^{23}$~W/cm$^2$) & $9\ee{16}$ & $2\ee{-45}$ & No & -- \\
SRF cavity (25~MV/m) & $3\ee{5}$ & $6\ee{-44}$ & No (direct) & SQMS (Q-ratio) \\
Solar active region & $4$ & $2\ee{-5}$ $^\dagger$ & Marginal & SDO/AIA \\
Magnetar surface & $4\ee{27}$ & $8\ee{-13}$ & Possible & Athena/XRISM \\
Neutron star merger & $\sim 10^{30}$ & $\sim 10^{-9}$ & Yes (GW?) & LIGO/ET \\
\bottomrule
\end{tabular}
\end{center}
$^\dagger$ Solar value uses coherent path integration over $\sim R_\odot$.


%=============================================================================
%=============================================================================
\part{Experiment 5: Vacuum EM Properties\\
in Strong Gravitational Fields}
\label{part:vacuum-em}
%=============================================================================
%=============================================================================


%=============================================================================
\section{The Proposal}
\label{sec:vacuum-em-proposal}
%=============================================================================

The vacuum loading picture makes a sharp prediction: the vacuum
permittivity $\epsilon_0$ and permeability $\mu_0$ are not fixed
constants but vary with the local $\psi$-field:
\begin{align}
\epsilon(\psi) &= \epsilon_0\, e^{(1+\kappa)\psi}, \\
\mu(\psi) &= \mu_0\, e^{(1-\kappa)\psi}.
\end{align}

Near a massive body with $\psi = -2\Phi/c^2$:
\begin{align}
\frac{\delta\epsilon}{\epsilon_0} &= (1+\kappa)\psi
\approx -(1+\kappa)\frac{2GM}{rc^2}, \\
\frac{\delta\mu}{\mu_0} &= (1-\kappa)\psi
\approx -(1-\kappa)\frac{2GM}{rc^2}.
\end{align}

For $\kappa = 0$ (symmetric loading): both change by $\psi$, and
the impedance $Z = \sqrt{\mu/\epsilon}$ is unchanged.  For
$\kappa \neq 0$: the impedance changes, producing a detectable
polarization-dependent effect.

\subsection{Has This Been Measured?}

\textbf{No.}  After extensive literature search, we find no experiment
that has directly measured $\epsilon_0$ or $\mu_0$ (or their ratio,
the vacuum impedance $Z_0 = \sqrt{\mu_0/\epsilon_0}$) as a function
of gravitational potential.

Related measurements that exist:
\begin{itemize}[nosep]
\item \textbf{Shapiro delay:}  Measures $n = c/v_{\text{ph}} = e^\psi$,
  which is the \textit{product} $\sqrt{\epsilon\mu}$, not the
  individual factors.

\item \textbf{Gravitational birefringence searches:}  Null results
  at $\Delta n \sim 10^{-28}$ (Ni et al., 1991).  These constrain
  the $\kappa$-dependent \textit{difference} $\epsilon - \mu$, but
  only through polarization rotation, not direct measurement.

\item \textbf{Gravity Probe B:}  Measured geodetic and frame-dragging
  precession of gyroscopes.  Not directly sensitive to $\epsilon_0$
  or $\mu_0$.

\item \textbf{Pound--Rebka:}  Measured gravitational redshift of
  photons.  Sensitive to $n(\psi)$ but not to the $\epsilon/\mu$
  split.
\end{itemize}

\textbf{The direct measurement of vacuum EM constitutive properties
in a gravitational field is a genuinely new experimental proposal.}

%=============================================================================
\section{Experimental Approaches}
\label{sec:vacuum-em-approaches}
%=============================================================================

\subsection{Approach A: Impedance Measurement via Reflection}

The vacuum impedance is $Z_0 = \sqrt{\mu_0/\epsilon_0}
\approx 376.730~\Omega$.  In a $\psi$-field:
\begin{equation}
Z(\psi) = Z_0\, e^{-\kappa\psi},
\end{equation}
so $\delta Z/Z_0 = -\kappa\psi$.

\paragraph{Signal at Earth's surface.}
$\psi_\oplus \approx -6.95\ee{-10}$ (due to Earth's gravity at
surface).  So:
\begin{equation}
\frac{\delta Z}{Z_0} = \kappa \times 6.95\ee{-10}.
\end{equation}
For $\kappa = 1$: $\delta Z/Z_0 \sim 7\ee{-10}$.

\paragraph{Measurement concept.}
An ultra-precise microwave reflectometry experiment measuring the
reflection coefficient of a precision impedance standard in a
varying gravitational potential (e.g., at different altitudes, or
using the annual solar potential modulation):
\begin{equation}
\delta\Gamma_{\text{refl}} \sim \frac{\delta Z}{Z_0}
\times \frac{\Delta\Phi}{c^2}.
\end{equation}

\paragraph{Sensitivity requirement.}
The annual solar modulation gives $\Delta\Phi/c^2 \sim 3.3\ee{-10}$.
Combined with $\kappa \sim 1$: $\delta Z/Z_0 \sim 3\ee{-10} \times
\kappa$.  Microwave reflectometry with network analyzers achieves
$\delta\Gamma \sim 10^{-6}$, a gap of 4 orders of magnitude.

Improvement path: SRF cavity techniques can measure impedance
changes at $\sim 10^{-10}$ through $Q$-factor shifts, potentially
closing the gap.

\subsection{Approach B: Polarization Rotation through $\nabla\psi$}

If $\kappa \neq 0$, a linearly polarized EM wave propagating through
a $\psi$-gradient acquires a polarization rotation:
\begin{equation}
\Delta\theta = \kappa \int \frac{d\psi}{ds}\, dl
= \kappa\,\Delta\psi.
\end{equation}

\paragraph{Solar limb measurement.}
For a radio signal grazing the Sun:
$\Delta\psi = 2\psi_\odot(R_\odot) \approx 8.4\ee{-6}$.
\begin{equation}
\Delta\theta = \kappa \times 8.4\ee{-6}~\text{rad}
\approx \kappa \times 0.48''~\text{(arcsec)}.
\end{equation}

For $\kappa = 1$: $\Delta\theta \approx 0.5''$.  This is a
\textit{large} signal by radio astronomy standards.

\paragraph{Complication: Faraday rotation.}
The solar corona produces Faraday rotation of order
$\Delta\theta_F \sim 10^3$~rad at 1~GHz (!) due to the dense
magnetized plasma.  The Faraday rotation scales as $f^{-2}$, while
the DFD vacuum polarization rotation is frequency-independent.

\textbf{Critical discriminant:}  Measure polarization rotation at
multiple frequencies.  The frequency-independent component is the
DFD signal.

\paragraph{Required instrument.}
A radio interferometer observing a linearly-polarized background
source (pulsar, AGN) during solar conjunction:
\begin{itemize}[nosep]
\item Frequency range: 1--100~GHz (to separate $f^{-2}$ Faraday
  from $f^0$ vacuum rotation)
\item Polarization purity: $< 10^{-4}$ instrumental polarization
\item Angular proximity: observations at $3$--$10\,\Rsun$ impact
  parameter
\item Multiple epochs: observe the same source at different solar
  elongations
\end{itemize}

\paragraph{Feasibility assessment.}
The VLA and ALMA routinely observe sources near the Sun for
Faraday rotation studies.  The multi-frequency analysis to extract
the frequency-independent component is technically straightforward.
The challenge is the small signal ($\sim 0.5'' \times \kappa$) amid
the large Faraday rotation.  However, the Faraday contribution is
well-modelled from density and magnetic field measurements.

\textbf{Estimated cost:}  Guest observer proposal at VLA + ALMA,
plus 1 postdoc for analysis.  Total: $\sim$\$150k.

\subsection{Approach C: TE/TM Cavity Frequency Splitting}

In an SRF cavity, TE and TM modes have different sensitivities to
$\epsilon$ and $\mu$.  A $\kappa \neq 0$ produces a differential
frequency shift between TE and TM modes:
\begin{equation}
\frac{\delta f_{\text{TE}} - \delta f_{\text{TM}}}{f}
\sim 2\kappa\psi.
\end{equation}

\paragraph{Signal using annual solar modulation.}
$\Delta\psi_\odot \sim 3.3\ee{-10}$, so:
\begin{equation}
\frac{\delta f_{\text{TE}} - \delta f_{\text{TM}}}{f}
\sim 2\kappa \times 3.3\ee{-10} \sim 7\ee{-10}\kappa.
\end{equation}

For $f = 1.3$~GHz: $\delta(f_{\text{TE}} - f_{\text{TM}}) \sim
0.9\kappa$~Hz.  With SQMS measurement precision of $\sim 10^{-3}$~Hz,
this gives sensitivity to $\kappa > 10^{-3}$.

\textbf{This is potentially the most sensitive laboratory approach
to measuring $\kappa$.}

The experiment piggybacks on the SQMS Phase~I Q-ratio measurement
(Part~\ref{part:srf}) at essentially zero incremental cost---one
simply adds TE/TM comparison to the measurement protocol.

%=============================================================================
\section{Summary and Ranking}
\label{sec:vacuum-em-summary}
%=============================================================================

\begin{center}
\begin{tabular}{lccl}
\toprule
\textbf{Approach} & \textbf{Sensitivity to $\kappa$}
& \textbf{Cost} & \textbf{Timeline} \\
\midrule
A: Impedance reflectometry & $\kappa > 10^3$ & \$500k & 2 yr \\
B: Solar polarization rotation & $\kappa > 1$ & \$150k & 1 yr \\
C: TE/TM cavity splitting (SQMS) & $\kappa > 10^{-3}$ &
  \$0 (incremental) & Immediate \\
\bottomrule
\end{tabular}
\end{center}

\textbf{Recommendation:}  Approach C (TE/TM splitting at SQMS) is
the highest-priority experiment for $\kappa$, as it is essentially
free if the Phase~I Q-ratio test proceeds.


%=============================================================================
%=============================================================================
\part{Experimental Roadmap: Integrated Timeline\\
and Priority Ranking}
\label{part:roadmap}
%=============================================================================
%=============================================================================


%=============================================================================
\section{Priority Ranking}
\label{sec:ranking}
%=============================================================================

We rank the five experimental proposals by a combined score of
scientific decisiveness, feasibility, and cost:

\begin{center}
\begin{tabular}{clcccl}
\toprule
\textbf{Rank} & \textbf{Experiment} & \textbf{Decisive?}
& \textbf{Feasible?} & \textbf{Cost} & \textbf{Critical Path} \\
\midrule
1 & SRF Q-ratio (SQMS) & Yes & Yes & \$3.2M & Fastest to decisive result \\
1a & Dark SRF add-on & Yes & Yes & \$450k & Fastest to any result \\
2 & TE/TM $\kappa$ (SQMS) & Yes & Yes & \$0* & Piggyback on \#1 \\
3 & Clock $\kappa$ (Th-229) & Yes & Developing & \$3--6M & Waiting on clock tech \\
4 & Corona archival & Partial & Yes & \$100k & Immediate start \\
5 & Solar polarization & If $\kappa\sim 1$ & Yes & \$150k & 1 yr \\
\bottomrule
\end{tabular}
\end{center}
*Incremental cost on SQMS Phase~I.

%=============================================================================
\section{Integrated Timeline}
\label{sec:timeline}
%=============================================================================

\subsection{Phase 0: Immediate Actions (0--12 months, \$250k)}

\begin{enumerate}
\item \textbf{SOHO/UVCS archival re-analysis:}  Hire a postdoc with
solar physics expertise.  Reanalyze O~VI double-transit data for
$\Gamma$.  Target: $\sigma_\Gamma < 0.5$.

\item \textbf{Dark SRF add-on proposal:}  Submit parasitic measurement
proposal to SQMS.  Add 2.4~GHz cavity to existing Dark SRF setup.
First multi-frequency Q-ratio data within 12 months.

\item \textbf{VLA/ALMA proposal:}  Submit guest observer proposal for
solar conjunction polarization observations.  Zero hardware cost;
telescope time is competitively allocated.
\end{enumerate}

\subsection{Phase I: SRF Campaign (12--36 months, \$3.2M)}

\begin{enumerate}
\item \textbf{SQMS Phase~I Q-ratio measurement:}  Full 4-cavity
programme with systematic controls.  Decisive binary test:
$\mathcal{R} = 0.49$ vs.\ 3.41.

\item \textbf{TE/TM splitting measurement:}  Simultaneous with
Q-ratio.  Sensitivity to $\kappa > 10^{-3}$.

\item \textbf{$\lambda$ bound improvement:}  From $3\ee{-5}$ to
$7.8\ee{-10}$ (null) or detection (positive).
\end{enumerate}

\subsection{Phase II: Clock Programme (36--72 months, \$3--6M)}

\begin{enumerate}
\item \textbf{Co-locate $^{229}$Th and $^{87}$Sr clocks:}
Leverage PTB/NIST nuclear clock development programmes.

\item \textbf{Begin 2-year measurement campaign:}  Annual solar
modulation of the nuclear-to-optical frequency ratio.

\item \textbf{Target:}  Measure or bound the ratio
$\delta\nu_{\text{nuc}}/\delta\nu_{\text{opt}} = 3000 \pm 500$.
\end{enumerate}

\subsection{Phase III: Next-Generation (72+ months)}

\begin{enumerate}
\item \textbf{Next-gen optical clocks ($10^{-21}$ stability):}
Sensitivity to $\kappa > 10^{-3}$ even with screening.

\item \textbf{Space-based clock comparison (ACES-2 class):}
Reduced screening in orbit; factor $\sim 40$ signal enhancement.

\item \textbf{Dedicated UV coronagraph (SMEX-class):}
Definitive $\Gamma$ measurement to $\pm 0.1$.

\item \textbf{Magnetar spectroscopy (Athena/NewAthena):}
EM-sourced $\psi$ perturbations from $B \sim 10^{11}$~T.
\end{enumerate}

%=============================================================================
\section{Total Programme Cost Summary}
\label{sec:total-cost}
%=============================================================================

\begin{center}
\begin{tabular}{lccl}
\toprule
\textbf{Phase} & \textbf{Cost} & \textbf{Timeline}
& \textbf{Key Deliverable} \\
\midrule
Phase 0 (immediate) & \$250k & 0--12 mo &
  First archival $\Gamma$; Dark SRF add-on \\
Phase I (SRF) & \$3.2M & 12--36 mo &
  Decisive Q-ratio; $\lbone < 10^{-9}$; $\kappa < 10^{-3}$ \\
Phase II (Clocks) & \$3--6M & 36--72 mo &
  Th/Sr ratio test; $\kappa$ from annual modulation \\
Phase III (next-gen) & \$150--300M & 72+ mo &
  Definitive $\Gamma$; space clocks; magnetar obs. \\
\midrule
\textbf{Phases 0--II total} & \textbf{\$6.5--9.5M} &
  \textbf{6 years} & \textbf{3 decisive tests} \\
\bottomrule
\end{tabular}
\end{center}


%=============================================================================
\section{Decision Tree}
\label{sec:decision-tree}
%=============================================================================

The experimental programme follows a decision tree based on
Phase~I results:

\paragraph{If Q-ratio $= 0.49 \pm 0.05$:}
DFD vacuum loading confirmed at $> 5\sigma$.  Proceed to clock
programme with high priority.  Apply for SQMS Phase~II
($\lambda$-measurement, \$15.8M).  Publish in Nature/Science.

\paragraph{If Q-ratio $= 3.41 \pm 0.34$:}
BCS confirmed; DFD vacuum loading Q-ratio prediction falsified.
Publish the null result (scientifically valuable as the world's
best bound on $\lbone$).  Reassess: the clock programme still
tests the $\kappa$-sector independently.

\paragraph{If Q-ratio intermediate ($1.0 < \mathcal{R} < 2.5$):}
Neither prediction confirmed.  Indicates either new physics beyond
both DFD and BCS, or uncontrolled systematics.  Design Phase~Ib
with enhanced systematic controls.

\paragraph{If archival $\Gamma = 4.0 \pm 0.3$:}
Strong support for the double-transit prediction.  Elevate solar
observation programme priority.  Propose Tier~2 (Solar Orbiter)
and Tier~3 (dedicated mission).

\paragraph{If archival $\Gamma \neq 4$ at $> 3\sigma$:}
Double-transit prediction falsified.  This constrains the vacuum
loading interpretation but does not affect the SRF or clock tests,
which probe different sectors of the theory.


%=============================================================================
\section{Conclusions}
\label{sec:conclusions}
%=============================================================================

We have developed the three falsifiable predictions of the vacuum
loading paper into five concrete experimental proposals with full
specifications:

\begin{enumerate}
\item \textbf{SRF Q-ratio test (SQMS):}  The single most
cost-effective experiment.  Decisive binary test between DFD
($\mathcal{R} = 0.49$) and BCS ($\mathcal{R} = 3.41$).  Can
leverage existing SQMS/Dark SRF infrastructure.  Cost: \$450k
(add-on) to \$3.2M (full programme).

\item \textbf{Solar corona double-transit:}  Archival SOHO/UVCS
re-analysis at \$100k can test $\Gamma = 4$ to $\sigma \sim 0.5$.
Definitive test requires a next-generation UV coronagraph.

\item \textbf{Clock comparison for $\kappa$:}  The $^{229}$Th
nuclear clock is the champion instrument, with a screening-independent
ratio prediction of 3000.  Current optical clocks provide
complementary data.  Timeline depends on nuclear clock maturation
(2027--2030).

\item \textbf{Strong EM field $\psi$ perturbation:}  Laboratory
fields produce undetectable $\psi$ perturbations.  Magnetar
observations with next-generation X-ray telescopes (Athena) are
the most promising route.

\item \textbf{Vacuum EM properties in gravity:}  No prior experiment
has measured $\epsilon_0$ or $\mu_0$ in a gravitational field.  The
TE/TM cavity frequency splitting at SQMS provides sensitivity to
$\kappa > 10^{-3}$ at zero incremental cost.  Solar polarization
rotation observations provide a complementary test.
\end{enumerate}

The critical path to first detection runs through the SQMS Q-ratio
test, with the nuclear clock programme providing an independent
confirmation channel on a 3--5 year timescale.  The total cost for
three decisive tests (Phases 0--II) is \$6.5--9.5M over 6 years.


%=============================================================================
\begin{thebibliography}{99}

\bibitem{Alcock2025Loading}
G.~T.~Alcock,
``Gravity as Vacuum Electromagnetic Loading,''
DFD Research Programme (2026).

\bibitem{Alcock2025DFD}
G.~T.~Alcock,
``Density Field Dynamics: A Complete Unified Theory,''
(2025); v32.

\bibitem{Alcock2025AbInitio}
G.~T.~Alcock,
``Ab Initio Derivation of the Fine Structure Constant
from Density Field Dynamics,''
(2025); v2.1.

\bibitem{Alcock2025Clock}
G.~T.~Alcock,
``Pre-Registration Protocol: Channel-Resolved Atomic Clock Tests
of Scalar-Field Gravitational Coupling in Density Field Dynamics,''
DFD Research Programme (2026); v4.

\bibitem{Alcock2025SQMS}
G.~T.~Alcock et al.,
``Precision Multi-Frequency Quality Factor Measurements in SRF
Cavities as a Test of Scalar-Field Modifications to Electrodynamics,''
SQMS Phase~I Proposal (2026).

\bibitem{Alcock2025EMCoupling}
G.~T.~Alcock,
``Deep Analysis of Electromagnetic Coupling to Scalar Refractive Fields,''
DFD Research Programme (2026).

\bibitem{Romanenko2023}
A.~Romanenko et al.,
``Dark SRF: first results from a dark photon search with
superconducting radio frequency cavities,''
\textit{Phys.\ Rev.\ Lett.} \textbf{131}, 261801 (2023).

\bibitem{Newkirk1961}
G.~Newkirk,
``The Solar Corona in Active Regions and the Thermal Origin of the
Slowly Varying Component of Solar Radio Radiation,''
\textit{Astrophys.\ J.} \textbf{133}, 983 (1961).

\bibitem{Ni1991}
W.-T.~Ni,
``Equivalence principles and precision experiments,''
in \textit{Precision Measurement and Fundamental Constants~II},
NBS Special Publication 617 (1991).

\bibitem{Peik2021}
E.~Peik et al.,
``Nuclear clocks for testing fundamental physics,''
\textit{Quantum Sci.\ Technol.} \textbf{6}, 034002 (2021).

\bibitem{Brewer2019}
S.~M.~Brewer et al.,
``$^{27}$Al$^+$ Quantum-Logic Clock with a Systematic Uncertainty
below $10^{-18}$,''
\textit{Phys.\ Rev.\ Lett.} \textbf{123}, 033201 (2019).

\bibitem{Bothwell2022}
T.~Bothwell et al.,
``Resolving the gravitational redshift across a millimetre-scale
atomic sample,''
\textit{Nature} \textbf{602}, 420 (2022).

\bibitem{Kohl1998}
J.~L.~Kohl et al.,
``UVCS/SOHO Empirical Determinations of Anisotropic Velocity
Distributions in the Solar Corona,''
\textit{Astrophys.\ J.\ Lett.} \textbf{501}, L127 (1998).

\bibitem{Guhathakurta1999}
M.~Guhathakurta, R.~R.~Fisher, and R.~C.~Altrock,
``Large-scale coronal temperature and density structures,''
\textit{Astrophys.\ J.} \textbf{521}, 441 (1999).

\end{thebibliography}

\end{document}
